% =============================================================
%  前言
% =============================================================
\Chapter{前言}

本书是为方便小豆角家庭辅导而作。平时工作忙，没有专门的时间来备课，
所以本书从编辑的第一天起就守一条原则：\textbf{编辑即备课}。

依据小豆角的\textbf{行为习惯、个人喜好、知识功底}，旨在去教材味，去教辅味，让孩子可无痛
、无感沉浸其中的方式加强数学思维，并把它做成一本只要拿在手上
\textbf{「照着念，对着做」}的执行手册，脑子在编辑的时候已经废完了，
后面只需要无脑对着做就行了。


\vspace{8pt}
\begin{warnbox}
\textbf{\tikz[baseline=-0.35ex]{\fill[warnc] (0,0.14) -- (0.26,-0.16) -- (-0.26,-0.16) -- cycle; \node[white,font=\fontsize{4.2}{4.2}\selectfont] at (0,-0.045) {!};}\ 注意}
\begin{enumerate}[label=\arabic*.,leftmargin=1.6em,itemsep=2pt,topsep=4pt]
  \item 作者不对此手册教学辅导效果作任何承诺。用本教材教乱了、教废了、没效果均为读者个人行为，与本书作者无关。
  \item 本书需配合北京师范大学出版社《数学》二年级上册（2025 年 7 月第 1 版，2022 课标修订）使用，其他教材未经验证，需谨慎评估。
\end{enumerate}
\end{warnbox}

\vspace{10pt}
本书采用 \LaTeX{} 编辑，\LaTeX{} 全书源码已全文开源，其他老师、家长可基于该开源项目
进行完善、优化或改编。

\vspace{4pt}

\begin{tipbox}\small
\textbf{开源地址}\\[4pt]
GitHub：\url{https://github.com/BD7IWD/20-Minute-Math-Quest}\\[3pt]
GitCode：\url{https://gitcode.com/BD7IWD/20-Minute-Math-Quest}

\vspace{6pt}
\textbf{开源说明}\\[4pt]
本书采用 \textbf{CC BY-NC-SA} 协议。\\
一句话概括：你可以随便改，改完后二次发布需要\textbf{按照}当前模式开源出来，
同时还不能拿去商用。
\end{tipbox}
